% ==============================================================================
% Social Science Paper Template  (dual-track: quantitative / qualitative)
%
% APA Publication Manual (7th ed.) for quantitative; Braun & Clarke (2006)
% thematic analysis for qualitative; Lincoln & Guba (1985) trustworthiness.
% Citation style: APA 7 author-year, approximated by natbib + apalike.bst.
%
% >>> PICK ONE TRACK <<<
%   QUANTITATIVE: keep Method / Results / Discussion sections as-is.
%   QUALITATIVE:  replace Method -> Methodology, Results -> Findings; expand
%                 the interview-transcript appendix and reflexivity block.
%
% Compile:  pdflatex main -> bibtex main -> pdflatex main -> pdflatex main
% ==============================================================================
\documentclass[11pt,a4paper]{article}

% ---- Page geometry ----
\usepackage[a4paper,top=25mm,bottom=25mm,left=25mm,right=25mm]{geometry}

% ---- Fonts & encoding (APA 7 recommends Times New Roman 11pt or Georgia 11pt) ----
\usepackage[T1]{fontenc}
\usepackage[utf8]{inputenc}
\usepackage{mathptmx}
\usepackage{helvet}
\usepackage{courier}

% ---- Line numbers + double spacing for review ----
\usepackage{lineno}
\usepackage{setspace}
\doublespacing                    % APA 7 manuscript: double-spaced
\linenumbers

% ---- Math, figures, tables ----
\usepackage{amsmath,amssymb}
\usepackage{graphicx}
\usepackage{booktabs}
\usepackage{array}
\usepackage{xcolor}
\usepackage{longtable}            % interview transcript tables
\usepackage{multirow}

% ---- Citations: APA-style author-year (approximated) ----
\usepackage[round,authoryear]{natbib}
\bibliographystyle{apalike}       % swap with plainnat / asa / chicago for sub-disciplines

% ---- Hyperref last ----
% \usepackage{hyperref}
% \hypersetup{colorlinks=true,linkcolor=blue,citecolor=blue,urlcolor=blue}

\title{A Concise and Informative Title of the Social Science Submission}
\author{First Author\textsuperscript{1,*}, Second Author\textsuperscript{2}, Senior Author\textsuperscript{1}\\
\textsuperscript{1}Department of X, University of Y, City, Country\\
\textsuperscript{2}Department of Z, Institute of W, City, Country\\
\textsuperscript{*}e-mail: first.author@univ.edu}
\date{}

\begin{document}
\maketitle

% ==============================================================================
% Abstract  (150-250 words; APA 7)
% ==============================================================================
\begin{abstract}
% TODO: 150-250 words. State the research question, method (paradigm +
% sample size + analysis), key findings, and implications. For quantitative
% papers list effect sizes and 95\% CIs.
\end{abstract}

\noindent\textbf{Keywords:} % TODO: 3-5 keywords separated by commas.
keyword one, keyword two, keyword three

% ==============================================================================
% 1. Introduction  (no "Introduction" heading per APA 7; CARS model)
% ==============================================================================
\section{Introduction}
\label{sec:intro}
% TODO: review prior literature; identify the gap; state the research
% question or hypotheses. For qualitative papers also include a brief
% research-paradigm statement (\texttt{\% TODO: research-paradigm}).

% ==============================================================================
% 2. Literature Review  (sometimes folded into Introduction; kept separate here)
% ==============================================================================
\section{Literature Review}
\label{sec:litreview}
% TODO: synthesise prior work; build the conceptual framework; for
% quantitative papers derive the testable hypotheses here (\texttt{\% TODO:
% hypothesis}).

% ==============================================================================
% 3. Methodology / Method
% ==============================================================================
% >>> QUANTITATIVE TRACK: section name = Method
% >>> QUALITATIVE TRACK:  section name = Methodology
\section{Methodology}
\label{sec:method}

\subsection{Research paradigm} % QUALITATIVE only
% TODO research-paradigm: state constructivist / interpretivist / critical
% stance and the methodological congruence rationale.

\subsection{Participants}
% TODO participants: sampling strategy (purposive / snowball / stratified
% random), inclusion criteria, sample size justification (saturation for
% qualitative; power analysis for quantitative), demographics table.

\subsection{Measures} % QUANTITATIVE only
% TODO measures: list each scale with its source, number of items, scoring,
% and reported reliability/validity evidence; report Cronbach's $\alpha$ in
% the present sample.

\subsection{Data collection}
% TODO procedure: interview protocol / focus-group guide / survey
% administration details; setting; language; duration; recording method.

\subsection{Data analysis}
% TODO analysis-quant: descriptive statistics, assumption checks, primary
% model (ANOVA / regression / SEM), software + version.
% TODO analysis-qual: Braun \& Clarke (2006) six-phase thematic analysis OR
% grounded-theory open/axial/selective coding; coder agreement; software
% (NVivo / ATLAS.ti / MAXQDA).

\subsection{Trustworthiness / Validity}
% QUALITATIVE: TODO: address Lincoln \& Guba (1985) credibility,
% transferability, dependability, confirmability.
% QUANTITATIVE: TODO: report construct validity (CFA), internal consistency,
% test-retest reliability.

\subsection{Ethical considerations}
% TODO irb: IRB approval number, informed consent procedure, confidentiality
% safeguards, protection of vulnerable populations, and any compensation.

% ==============================================================================
% 4. Findings / Results
% ==============================================================================
% >>> QUANTITATIVE TRACK: section name = Results
% >>> QUALITATIVE TRACK:  section name = Findings
\section{Findings}
\label{sec:findings}

% --- QUALITATIVE: thematic tree presentation ---
\subsection{Theme 1: a short, descriptive theme title}
% TODO themes: open with one paragraph summarising the theme, then 2-3
% indented participant quotations with identifiers, e.g.
% \begin{quote}``...'' (P7, female, age 34)\end{quote}
% followed by interpretive commentary linking back to the literature.

\subsection{Theme 2: another descriptive theme title}
% TODO: as above.

% --- QUANTITATIVE: descriptive + inferential ---
% \subsection{Descriptive statistics}
% \subsection{Primary analysis}
% \subsection{Supplementary analyses}

\begin{table}[t]
\caption{Participant demographics. Caption ABOVE the table (APA convention).}
\label{tab:demographics}
\centering
\begin{tabular}{lcc}
\toprule
Characteristic & Category & $n$ (\%) \\
\midrule
Gender         & Female   & --- (---) \\
               & Male     & --- (---) \\
Age (years)    & $M \pm SD$ & --- $\pm$ --- \\
\bottomrule
\end{tabular}
\end{table}

% ==============================================================================
% 5. Discussion
% ==============================================================================
\section{Discussion}
\label{sec:discussion}
% TODO: (i) restate the key findings without statistics, (ii) interpret in
% light of prior literature, (iii) theoretical contribution, (iv) practical
% / policy implications, (v) limitations, (vi) future research.

\subsection{Reflexivity} % QUALITATIVE only
% TODO reflexivity: discuss the researcher's positionality, prior
% assumptions, and how the researcher-participant relationship may have
% shaped the findings.

\subsection{Limitations}
% TODO: sampling bias, self-report bias, single-setting generalisability,
% cross-sectional design, etc.

% ==============================================================================
% 6. Conclusion
% ==============================================================================
\section{Conclusion}
\label{sec:conclusion}
% TODO: one paragraph on the contribution and a forward-looking statement.

% ==============================================================================
% Declarations
% ==============================================================================
\section*{Declarations}

\noindent\textbf{Funding:} % TODO: grant numbers and funding bodies.

\noindent\textbf{Conflicts of interest:} % TODO: declare or state none.

\noindent\textbf{Ethics approval:} % TODO irb: IRB name + approval number.

\noindent\textbf{Informed consent:} % TODO: state written / verbal consent procedure.

\noindent\textbf{Data availability:} % TODO: for quantitative, deposit data
% + codebook; for qualitative, state how transcripts are anonymised and
% whether they can be shared (often restricted to protect participants).

% ==============================================================================
% Acknowledgements
% ==============================================================================
\section*{Acknowledgements}
% TODO: funders, translators, transcribers, community partners, participants.

% ==============================================================================
% References  —  APA 7 author-year style
% Compile:  pdflatex main -> bibtex main -> pdflatex main -> pdflatex main
% ==============================================================================
\bibliography{references}
% Inline fallback (uncomment if not using BibTeX):
% \begin{thebibliography}{99}
% \bibitem{ref_example} Author, F., \& Author, S. (20YY). Title of the paper.
%   \emph{Journal Name, XX}(Y), 1--10.
% \end{thebibliography}

% ==============================================================================
% Appendices  —  interview guide, codebook, full transcripts (qualitative)
% ==============================================================================
\appendix

\section{Interview Guide}
\label{app:interview-guide}
% TODO: list the semi-structured interview questions and probes.

\section{Codebook}
\label{app:codebook}
% TODO qualitative: list themes, sub-themes, definitions, and exemplar quotes.

\section{Interview Transcripts}
\label{app:transcripts}
% TODO transcripts: provide de-identified transcripts formatted as longtable
% rows (timestamp | participant ID | utterance), or state the repository
% where the transcripts are deposited (with access criteria).

\end{document}
